\uuid{31N8}
\exo7id{7804}
\titre{exo7 7804}
\auteur{mourougane}
\organisation{exo7}
\datecreate{2021-08-11}
\isIndication{false}
\isCorrection{false}
\chapitre{Groupe orthogonal et quaternions}
\sousChapitre{Groupe orthogonal et quaternions}
\cpge{mpsi-mp2i.applications-lineaires.b, pcsi.applications-lineaires.b, ptsi.applications-lineaires.b, mpsi-mp2i.espaces-prehilbertiens.c, pcsi.espaces-prehilbertiens.c, mp-mpi.endomorphismes-euclidien.b, pc.endomorphismes-euclidiens.b, psi.endomorphismes-euclidien.b, pt.isometries-euclidien.b, mp-mpi.endomorphismes-euclidien.c, pc.endomorphismes-euclidiens.a, psi.endomorphismes-euclidien.a, pt.isometries-euclidien.a, mp-mpi.endomorphismes-euclidien.e, psi.endomorphismes-euclidien.e, pt.isometries-euclidien.d}
\module{Algèbre et géométrie}
\niveau{L3}
\difficulte{}

\contenu{
\texte{
Soit $f$ une forme bilinéaire symétrique non dégénérée sur un espace vectoriel $E$.
}
\begin{enumerate}
    \item \question{Soit $u\in GL(E)$ une involution. Montrer que $u$ est orthogonale si
et seulement si $E_+(u)=E_+:=Ker (u-Id)$ et $E_-:=Ker(u+Id)$ sont
orthogonaux.
 Montrer alors que $(E_+)^\perp = E_-$ et que $E_+$ n'est pas isotrope.}
    \item \question{Soit $F\subset E$ un sous-espace non isotrope. Montrer qu'il
 existe une unique involution orthogonale telle que $E_+(u)=F$.}
    \item \question{Montrer que $O(f)\equiv SO(f)\ltimes \{-1,1\}$.}
\end{enumerate}
}
